% ============================================================
\chapter{Identificación de requisitos}
\label{cap:requisitos}
% ============================================================

Este capítulo identifica y justifica los requisitos que condicionaron el diseño
de \textit{kubescan}. Los requisitos funcionales y no funcionales se derivan del
análisis de vectores de ataque documentado en el estado del arte y de las
limitaciones detectadas en las herramientas existentes. Delimitan además el
alcance del sistema: qué analiza, sobre qué capas se apoya y qué queda fuera de
sus objetivos. El capítulo cierra con dos estudios previos que determinaron
directamente el diseño de la herramienta: la estrategia de partición y
validación de datos, y las métricas empleadas para evaluarla.

% -----------------------------------------------------------
\section{Requisitos del sistema}
\label{sec:requisitos}
% -----------------------------------------------------------

El diseño de \textit{kubescan} parte de un análisis de requisitos derivado tanto de
las limitaciones identificadas en el estado del arte (Sección~\ref{sec:sintesis}) como
de los vectores de ataque documentados por la \ac{NSA} y la \ac{CISA} \citep{NSA2022}.
Este análisis se traduce en dos conjuntos de requisitos: funcionales, que definen
las capacidades que el sistema debe ofrecer, y no funcionales, que acotan sus
propiedades de calidad. Ambos conjuntos se detallan a continuación, junto con la
justificación concreta de cada uno.

\subsection{Requisitos funcionales}

\begin{description}
  \item[RF-1 — Detección de misconfiguraciones individuales.] El sistema debe cumplir
    dos obligaciones separables sobre cada manifiesto \acs{YAML}. La primera es
    clasificarlo como seguro o \textit{misconfigured}. La segunda es producir un
    índice de severidad en tres niveles (limpio, bajo\,--\,medio, alto\,--\,crítico)
    que ordene las misconfiguraciones detectadas por gravedad.
    \textit{Criterio de aceptación:} ambas obligaciones se consideran satisfechas
    cuando cada una de las dos clasificaciones alcanza F1 macro superior a 0,85
    sobre el conjunto de test reservado, umbral adoptado de OE2
    (Capítulo~\ref{cap:objetivos}) y aplicado también al índice de severidad por
    tratarse del mismo clasificador sobre el mismo corpus.
    \textit{Justificación:} las herramientas de análisis estático existentes
    (\textit{Checkov}, \textit{kube-linter}) ya cubren esta capacidad a nivel de
    manifiesto individual (Sección~\ref{sec:iac}); el sistema debe igualarla como
    condición previa a razonar sobre la estructura del clúster, y el índice de
    severidad aporta priorización que las herramientas de reglas fijas no ofrecen.
  \item[RF-2 — Detección de cadenas de ataque.] El sistema debe clasificar un conjunto
    de manifiestos que describe un clúster en una de tres categorías: \textit{clean},
    \textit{isolated} (misconfiguraciones sin cadena viable) o
    \textit{attack\_chain} (ruta de escalada de privilegios hasta movimiento lateral).
    \textit{Criterio de aceptación:} toda ejecución debe emitir exactamente una de
    esas tres categorías por clúster analizado, expuesta en la salida \acs{JSON}, y
    las cadenas de ataque del conjunto de test reservado deben quedar situadas en
    las primeras posiciones del \textit{ranking} de riesgo con arreglo al umbral
    fijado en RF-5. La proporción de cadenas clasificadas correctamente por
    \textit{argmax} se reporta como métrica descriptiva del comportamiento del
    modelo y no constituye por sí sola condición de aceptación, dado que el caso
    de uso operativo del sistema es la priorización y no el etiquetado exacto.
    \textit{Justificación:} es la carencia central identificada en el estado del arte
    (Sección~\ref{sec:sintesis}): ninguna herramienta de análisis estático por
    manifiesto individual emplea un modelo supervisado que razone sobre las
    relaciones entre recursos que habilitan una cadena, y las comprobaciones de
    grafo basadas en reglas declarativas que incorpora \textit{Checkov} se limitan
    a dependencias conocidas de antemano.
  \item[RF-3 — Interfaz de línea de comandos.] El sistema debe exponer un comando
    \texttt{kubescan scan <dir>} para análisis de directorios locales y un comando
    \texttt{kubescan live} para análisis de clústeres activos mediante \textit{kubectl}.
    \textit{Justificación:} un control de seguridad solo es adoptable en la práctica
    si se integra en los flujos de trabajo existentes de un equipo de seguridad o de
    \acs{CICD}; una interfaz de línea de comandos es el punto de integración mínimo
    común a ambos.
  \item[RF-4 — Formatos de salida.] El sistema debe producir informes en formato texto
    legible y en formato \acs{JSON} estructurado para integración con herramientas externas.
    \textit{Justificación:} el formato texto sirve a la revisión manual de un
    operador, mientras que el formato \acs{JSON} permite el consumo automatizado
    del veredicto por otras herramientas del \textit{pipeline} de \acs{CICD}.
  \item[RF-5 — Priorización de riesgos.] El sistema debe producir una clasificación
    ordenada de clústeres por puntuación de riesgo, de modo que las cadenas de
    ataque reales queden situadas en las primeras posiciones.
    \textit{Criterio de aceptación:} Precisión@5 superior a 0,70 sobre el
    \textit{ranking} del sistema desplegado en el conjunto de test reservado.
    \textit{Alcance del umbral:} la cifra de 0,70 se adopta de OE4
    (Capítulo~\ref{cap:objetivos}), donde se define como media de validación
    cruzada de 5~pliegues de la Capa~2; el presente requisito la traslada al
    \textit{ranking} del sistema completo, de manera que ambas medidas comparten
    umbral pero no protocolo de cálculo.
    \textit{Precisión estadística:} el tamaño del conjunto de test reservado
    (86~grafos, 5~cadenas) hace que el intervalo de confianza \textit{bootstrap}
    al 95\,\% de la Precisión@5 abarque \([0{,}60,\ 1{,}00]\) y por tanto contenga
    el propio umbral, de modo que la superación del criterio debe leerse como
    consistente con el objetivo de diseño y no como una garantía estadística.
    \textit{Justificación:} un equipo de seguridad revisa manualmente un número
    reducido de alertas por jornada; un ranking que sitúe las cadenas reales en
    las primeras posiciones hace viable esa revisión sobre un volumen de
    clústeres que de otro modo sería inabarcable.
\end{description}

\subsection{Requisitos no funcionales}

\begin{description}
  \item[RNF-1 — Trazabilidad.] Todas las reglas de detección de nodos de escape y
    movimiento lateral deben ser explícitas, auditables y derivadas de fuentes
    normativas documentadas \citep{NSA2022,BishopFox2021}.
    \textit{Criterio de aceptación:} cada veredicto debe poder rastrearse a las
    banderas concretas que lo originan, expuestas en la salida \acs{JSON} junto a
    la identificación de los nodos con capacidad de escape, y cada regla debe
    remitir a una fuente normativa citable.
    \textit{Justificación:} un sistema que informa decisiones de seguridad debe
    permitir que su veredicto sea auditado y cuestionado; basar cada regla en una
    fuente pública evita además introducir vectores de ataque no documentados.
  \item[RNF-2 — Reproducibilidad.] Los modelos entrenados deben almacenarse como
    \textit{checkpoints} que el paquete pueda localizar sin intervención del
    usuario y que se distribuyan junto al repositorio del proyecto, de forma que
    el análisis sea idempotente dado el mismo directorio de entrada.
    \textit{Criterio de aceptación:} el requisito se descompone en dos condiciones
    comprobables por separado. La primera es de disponibilidad: los pesos deben
    resolverse mediante la cadena documentada de resolución de \textit{checkpoints}
    —opción explícita de línea de comandos, variable de entorno y ubicación por
    defecto dentro del paquete— y quedar publicados en el repositorio distribuido,
    sin descarga adicional por parte del usuario. La segunda es de determinismo:
    dos ejecuciones consecutivas sobre el mismo directorio de manifiestos y el
    mismo conjunto de pesos deben producir informes \acs{JSON} idénticos.
    \textit{Justificación:} un control de seguridad con salidas no deterministas
    no es fiable como \textit{gate} de \acs{CICD}, donde el mismo manifiesto debe
    producir siempre el mismo veredicto.
  \item[RNF-3 — Instalabilidad.] El sistema debe ser instalable mediante
    \texttt{pip install -e kubescan/} sin dependencias de compiladores externos.
    \textit{Criterio de aceptación:} la instalación debe completarse sin error en
    intérpretes Python 3.10 o superior, cota declarada en OE6
    (Capítulo~\ref{cap:objetivos}) y en los metadatos de empaquetado, resolviendo
    todas las dependencias sin invocar un compilador del sistema.
    \textit{Justificación:} reduce la fricción de adopción por un equipo de
    seguridad que no necesariamente dispone de un entorno de compilación
    preparado para dependencias nativas.
  \item[RNF-4 — Latencia.] El análisis de un clúster de hasta 500 manifiestos debe
    completarse en menos de 60 segundos en hardware de gama media (CPU).
    \textit{Criterio de aceptación:} la cota de 500~manifiestos se formula como
    límite de diseño y no como tamaño de un caso de prueba disponible, ya que el
    mayor clúster del corpus recopilado no alcanza esa dimensión. La verificación
    consiste, por tanto, en medir el tiempo de análisis sobre clústeres de tamaño
    creciente y comprobar que el coste por manifiesto observado mantiene el tiempo
    proyectado a 500~manifiestos por debajo del límite de 60~segundos.
    \textit{Justificación:} un tiempo de análisis del orden de segundos es
    condición necesaria para que el sistema pueda operar como control preventivo
    dentro de un \textit{pipeline} de \acs{CICD} sin introducir un cuello de
    botella perceptible.
\end{description}

\subsection{Trazabilidad de los requisitos}

Cada requisito se materializa en una decisión de diseño concreta y se contrasta
después con una evidencia experimental u operacional determinada. La
Tabla~\ref{tab:trazabilidad} recoge esa correspondencia, indicando para cada
identificador la sección del Capítulo~\ref{cap:herramienta} que describe el
componente encargado de satisfacerlo y la sección del
Capítulo~\ref{cap:evaluacion} que aporta la evidencia de su cumplimiento. Los
requisitos que no admiten una métrica predictiva dedicada, como los de interfaz,
formatos de salida y propiedades no funcionales, se concentran en la
verificación operacional de la Sección~\ref{sec:verificacion_requisitos}.

\begin{table}[htbp]
\centering\small
\caption{Trazabilidad de requisitos}
\label{tab:trazabilidad}
\begin{tabular}{@{}lp{5.1cm}ll@{}}
\toprule
ID & Requisito & Diseño & Verificación \\
\midrule
RF-1  & Misconfiguraciones individuales & \ref{sec:capa1_diseno}  & \ref{sec:layer1} \\
RF-2  & Cadenas de ataque               & \ref{sec:capa2_diseno}  & \ref{sec:layer2} \\
RF-3  & Interfaz de línea de comandos   & \ref{sec:implementacion} & \ref{sec:verificacion_requisitos} \\
RF-4  & Formatos de salida              & \ref{sec:implementacion} & \ref{sec:verificacion_requisitos} \\
RF-5  & Priorización de riesgos         & \ref{sec:capa3_diseno}  & \ref{sec:metrica_global} \\
RNF-1 & Trazabilidad                    & \ref{sec:features_table} & \ref{sec:verificacion_requisitos} \\
RNF-2 & Reproducibilidad                & \ref{sec:implementacion} & \ref{sec:verificacion_requisitos} \\
RNF-3 & Instalabilidad                  & \ref{sec:implementacion} & \ref{sec:verificacion_requisitos} \\
RNF-4 & Latencia                        & \ref{sec:implementacion} & \ref{sec:verificacion_requisitos} \\
\bottomrule
\end{tabular}
\fuentepropia
\end{table}

El índice de severidad exigido por RF-1 se verifica de forma diferenciada
respecto de la clasificación binaria, mediante los resultados por clase de la
Tabla~\ref{tab:severity_f1}. La restricción de viabilidad estructural que
condiciona el \textit{ranking} evaluado en RF-5 se describe en la
Sección~\ref{sec:gate_estructural} y su efecto se cuantifica en la
Sección~\ref{sec:layer3}.

% -----------------------------------------------------------
\section{Justificación del alcance}
\label{sec:alcance}
% -----------------------------------------------------------

El sistema opera principalmente sobre manifiestos \acs{YAML} estáticos, sin
requerir acceso a un clúster activo, aunque incorpora opcionalmente una vía de
consulta directa al clúster mediante \texttt{kubescan live}. La elección del
análisis estático como vía principal responde al principio \textit{shift-left
security}: detectar una cadena de ataque antes del despliegue evita el coste de
remediarla en producción, mientras que las herramientas de grafo de ataque sobre
clúster activo (\textit{KubeHound} \citep{KubeHound2023}, \textit{IceKube}
\citep{IceKube2023}) solo pueden operar \textit{a posteriori} del despliegue y
requieren acceso con privilegios elevados a \texttt{kubectl} sobre un clúster ya
en ejecución. La vía de consulta directa se mantiene
como capacidad complementaria para auditorías puntuales sobre clústeres ya
desplegados, no como sustituto de la vía preventiva.

La arquitectura de tres capas responde a una separación de responsabilidades
entre dos niveles de análisis y su combinación: la Capa~1 evalúa cada manifiesto
de forma aislada (RF-1), la Capa~2 razona sobre la estructura relacional del
clúster que ningún manifiesto individual expone por sí solo (RF-2), y la Capa~3
combina ambas señales en un único veredicto priorizable (RF-5). Ninguna capa por
sí sola satisface el conjunto completo de requisitos funcionales: un clasificador
de manifiesto individual no puede expresar una relación entre dos recursos, y un
modelo de grafo sin las características de nodo de la Capa~1 pierde la señal de
riesgo específica de cada manifiesto. La elección de \textit{Random Forest} para
la Capa~1 está además condicionada por tres propiedades que la
Tabla~\ref{tab:ml_comparison} contrasta entre algoritmos candidatos: su
independencia respecto de la normalización de escala de las entradas, su
tolerancia al desequilibrio de clases y la granularidad de la puntuación continua
que produce, requisito este último derivado de que dicha salida se emplea como
característica de nodo en la Capa~2.

A las tres capas se añade una cuarta componente arquitectónica: la restricción
de viabilidad estructural descrita en la Sección~\ref{sec:gate_estructural}, que
excluye del \textit{ranking} de riesgo aquellos clústeres cuya topología no
puede albergar una cadena de ataque, como los compuestos por un único recurso.
Su inclusión no responde a una preferencia de diseño sino a una necesidad
verificada empíricamente: sin ella, la Precisión@1 sobre el conjunto de test
reservado cae a 0,00 y la Precisión@5 a 0,60, valor este último situado por
debajo del umbral que el propio RF-5 fija como criterio de aceptación. La
restricción actúa así como condición previa de admisibilidad al \textit{ranking},
complementaria a las señales aprendidas por los tres modelos y no sustitutiva de
ninguna de ellas.

Quedan explícitamente fuera del alcance del sistema: la explotación activa de las
cadenas de ataque detectadas (el sistema identifica rutas potenciales, no las
verifica mediante ataque controlado); la remediación automática de las
misconfiguraciones encontradas; y el análisis de vulnerabilidades a nivel de
imagen de contenedor (\textit{CVE scanning}), cubierto por herramientas
especializadas de escaneo de imágenes y complementario, no sustitutivo, del
análisis estructural que aporta \textit{kubescan}.

% -----------------------------------------------------------
\section{Estrategia de partición y validación de datos}
\label{sec:particion}
% -----------------------------------------------------------

La estrategia de partición de datos difiere entre la Capa~1 (datos tabulares) y la
Capa~2 (datos de grafos), dado que sus tamaños y distribuciones son distintos. Esta
estrategia condiciona directamente la comprobabilidad de RF-1 y RF-2, cuyos
criterios de aceptación se formulan sobre un conjunto de test reservado, y de
RF-5, cuyo umbral de Precisión@5 se mide sobre el \textit{ranking} producido para
ese mismo conjunto. El tamaño final de cada partición depende del proceso de
adquisición y construcción descrito en las
Secciones~\ref{sec:construccion_dataset} y~\ref{sec:construccion_grafo} del
Capítulo~\ref{cap:herramienta}.

Para la Capa~1, el conjunto de manifiestos se divide en un conjunto de entrenamiento
(80\,\%) y un conjunto de test (20\,\%), con estratificación por la variable objetivo
para mantener la proporción de clases en ambos subconjuntos. La validación cruzada
de 5~pliegues aplicada sobre el corpus tabular constituye un protocolo de
estimación del rendimiento esperado, no un procedimiento de búsqueda de
hiperparámetros: la configuración del clasificador es fija y se documenta en la
Sección~\ref{sec:capa1_diseno}, sin que intervenga ninguna exploración de rejilla
ni muestreo aleatorio del espacio de configuraciones. Al no existir búsqueda de
hiperparámetros, ninguna decisión de modelado queda condicionada por el
rendimiento observado, lo que preserva la validez del conjunto de test como
medida final del criterio de aceptación de RF-1.

Ambos protocolos difieren, no obstante, en el corpus sobre el que operan, y esa
diferencia debe explicitarse para interpretar correctamente los resultados de la
Capa~1. La validación cruzada de 5~pliegues se ejecuta sobre la totalidad del
corpus tabular, es decir, sobre las mismas filas que después integran el conjunto
de test, mientras que la columna de test procede del modelo entrenado únicamente
sobre el 80\,\% de entrenamiento. En consecuencia, la cifra de validación cruzada
no constituye una estimación sobre datos completamente ajenos al ajuste y debe
leerse como indicador de estabilidad del rendimiento entre particiones del
corpus; la única estimación sobre datos retenidos es la de la columna de test,
que es la que sustenta la verificación de RF-1.

Para la Capa~2 se reserva desde el inicio una partición de test de clústeres
originales, y el resto se evalúa mediante validación cruzada estratificada de
5~pliegues, con estratificación por etiqueta de clúster (\textit{clean},
\textit{isolated}, \textit{attack\_chain}) y particionado consciente de grupos y
de familias de plantilla, cuyo procedimiento detalla la
Sección~\ref{sec:aumentacion}. La partición global reserva 815~grafos para
entrenamiento, 88 para validación y 86 para test, mientras que el conjunto sobre
el que operan los pliegues de validación cruzada comprende 513~clústeres
originales. La validación cruzada complementa al conjunto de test reservado porque
es la estrategia estándar para conjuntos de grafos pequeños \citep{Hamilton2020}:
con pocas muestras de la clase minoritaria, una única partición de validación no
sería representativa por sí sola, de modo que el criterio de aceptación de RF-2
requiere ambas evidencias.

% -----------------------------------------------------------
\section{Métricas de evaluación}
\label{sec:metricas}
% -----------------------------------------------------------

La selección de métricas de evaluación responde a los requisitos operacionales del
sistema y a las características del dataset. Cada métrica se corresponde con el
criterio de aceptación de un requisito concreto: la F1 macro cuantifica el
criterio de RF-1 y describe la clasificación por clases invocada en RF-2, la
Precisión@5 cuantifica el criterio de RF-5, y la \ac{AUC}-\acs{ROC} y la
\ac{FPR}\_clean caracterizan el margen operativo con el que ambos criterios se
satisfacen. Todas ellas se calculan sobre las particiones definidas en la
Sección~\ref{sec:particion}, cuyo alcance determina si el valor resultante
constituye una estimación sobre datos retenidos. Las cuatro se definen a
continuación junto con la razón por la que resultan preferibles a alternativas de
uso más extendido.

\paragraph{F1 macro.} Se aplica a la Capa~1 y a la clasificación por clases de la
Capa~2. La F1 macro calcula la media no ponderada de la F1 de cada clase, otorgando igual
importancia a las clases minoritarias (misconfiguraciones críticas) que a las
mayoritarias (manifiestos seguros). Esta propiedad es fundamental en datasets de
seguridad desequilibrados, donde la exactitud (\textit{accuracy}) sería una métrica
engañosa —sobre el conjunto de test de la Capa~1 (566~muestras), un clasificador
que predice siempre «seguro» obtendría una exactitud del 56,3\,\% pero F1 macro
del 36,0\,\%.

\paragraph{\ac{AUC}-\acs{ROC}.} Se aplica al modelo binario de la Capa~1 y mide
la capacidad del clasificador para
separar las dos clases independientemente del umbral de decisión, lo que permite
al operador ajustar la sensibilidad según las necesidades del entorno.

\paragraph{Precisión@k.} Se aplica a la Capa~2 y al \textit{ranking} del sistema
completo. En lugar de evaluar la clasificación por umbral, P@k cuenta los
clústeres de cadena de ataque situados en las primeras \(k\) posiciones del
\textit{ranking} por puntuación de riesgo y normaliza ese recuento por
\(\min(k,\, n_{\text{pos}})\), siendo \(n_{\text{pos}}\) el número de cadenas
presentes en el conjunto evaluado:

\begin{equation}
  \text{P@}k \;=\;
  \frac{\left|\left\{\, i \in \text{top-}k \;:\; y_i = \textit{attack\_chain} \,\right\}\right|}
       {\min(k,\, n_{\text{pos}})}
  \label{eq:precision_at_k}
\end{equation}

Esta definición coincide con la fracción convencional sobre \(k\) únicamente
cuando el conjunto evaluado contiene al menos \(k\) cadenas; cuando contiene
menos, el denominador pasa a ser \(n_{\text{pos}}\) y la métrica alcanza el valor
1 si y solo si todas las cadenas existentes ocupan posiciones dentro del top-\(k\).
La forma normalizada es la que dota de sentido al techo teórico de la métrica
analizado en la Sección~\ref{sec:layer2}: sin ella, un conjunto con menos de
\(k\) cadenas tendría un máximo alcanzable inferior a la unidad por construcción
y no por deficiencia del modelo. Esta métrica está directamente
alineada con el caso de uso operacional: un equipo de seguridad que audita los
\(k\) clústeres de mayor riesgo en su infraestructura quiere que todos sean cadenas
de ataque reales, no falsas alarmas. La elección de \(k=5\) refleja la capacidad
de revisión manual típica de un equipo de seguridad reducido en una jornada
laboral, valor de \(k\) que coincide con el fijado en el criterio de aceptación
de RF-5.

\paragraph{\ac{FPR}\_clean.} Esta tasa de falsos positivos sobre clústeres
limpios mide la fracción de clústeres limpios entre los \(k\) primeros del \textit{ranking} de
riesgo\footnote{Operacionalmente es una tasa de limpios en el top-\(k\), no una
tasa de falsos positivos en sentido estricto (que normalizaría por el total de
clústeres limpios); se mantiene la denominación \ac{FPR}\_clean por coherencia con
la implementación.}.
Un \ac{FPR}\_clean\,=\,0 significa que ningún clúster limpio aparece entre los de mayor
puntuación, eliminando las investigaciones innecesarias.
